\documentclass[11pt]{article}
\usepackage[margin=1in]{geometry}
\usepackage{array}
\usepackage{booktabs}
\usepackage{hyperref}
\usepackage{longtable}
\usepackage{graphicx}
\usepackage{xcolor}

\title{Variations of Monte Carlo Tree Search for a Minecraft Sword PvP Agent}
\author{Nguyen Huy An Khanh}
\date{COMP2050 Programming Project\\Instructor: Leandro S. Marcolino\\June 21, 2026}

\begin{document}

\maketitle

\begin{center}
\textbf{Project repository:} \url{https://github.com/Ankhanh975/Programing-Project-2}
\end{center}

\begin{abstract}
This project studies Monte Carlo Tree Search (MCTS) for real-time sword combat in a Minecraft-like environment. The agent is implemented in TypeScript with Mineflayer for server interaction and prismarine-physics for a repeatable local duel simulator. Beyond a standard MCTS controller, the project implements and evaluates several variations: hybrid MCTS with a heuristic action bonus, risk-aware MCTS, high-exploration MCTS, and higher-budget MCTS. The final experiment generator produced 42 duel runs and 4,598 logged decisions, including 40 repeated strategy trials and two stress-test scenarios. The results show that the planner learns plausible combat behavior, especially closing distance before attacking, but also that stronger search settings can matter more than a hand-designed heuristic bonus in this simulator.
\end{abstract}

\section{Introduction}

Real-time game combat is a useful setting for studying artificial intelligence because an agent must make decisions under time pressure while handling movement, positioning, and uncertainty about the opponent. Minecraft sword PvP is a good example: a competent fighter must approach, maintain useful spacing, wait for attack range, and avoid wasting attacks during cooldown.

The goal of this project is to build and evaluate a sword-duel agent. The base algorithm is Monte Carlo Tree Search, a planning method that estimates action values by simulating possible futures. The project then goes beyond standard MCTS by adding domain-specific variations and comparing their effects experimentally.

The research questions are:
\begin{enumerate}
  \item Can a small MCTS planner produce plausible Minecraft sword-duel behavior in a physics simulator?
  \item How do creative variations such as heuristic evaluation, risk-aware scoring, exploration control, and larger search budgets change the agent's behavior?
  \item Which measurements best explain the observed combat outcomes?
\end{enumerate}

\section{Related Work}

Monte Carlo Tree Search was popularized by the UCT algorithm introduced by Kocsis and Szepesvári, which combines tree search with upper-confidence exploration. MCTS is widely used in game AI because it can choose actions without exhaustively searching the full game tree. The algorithm repeats selection, expansion, simulation, and backpropagation while balancing exploitation and exploration.

A major milestone was the original AlphaGo system (Silver et al., 2016), which combined MCTS with deep policy and value networks to achieve superhuman performance in Go. AlphaGo demonstrated that tree search can be significantly strengthened when guided by learned models. Although this project does not use neural networks, AlphaGo motivates the idea of enhancing a standard MCTS framework with additional guidance mechanisms. The same planning idea can also be adapted to real-time games by limiting the number of search iterations per decision.

Minecraft bot development commonly uses Mineflayer, a JavaScript API that exposes player entities, inventory, movement controls, chat, and attack actions. Prismarine-physics approximates Minecraft movement and collision behavior locally. In this project, these libraries provide the environment layer, while the combat search logic is implemented separately.

The project differs from standard turn-based examples because the action space is built from Minecraft controls: moving forward, sprinting, jumping, strafing, retreating, and attacking. MCTS therefore searches over short continuous-position futures generated by a physics simulator, while still choosing from a discrete action set.

\section{Methodology}

\subsection{Environment}

The implementation is written in TypeScript. The main runtime entry point is \texttt{src/bot.ts}; shared combat logic is in \texttt{src/combat.ts}; and report data generation is in \texttt{src/report-materials.ts}. The project supports two modes:
\begin{itemize}
  \item \textbf{Server mode}: a Mineflayer bot connects to a Minecraft server, equips a sword, identifies an opponent, and applies the selected movement controls and attacks.
  \item \textbf{Physics mode}: two local simulated fighters duel in a flat prismarine-physics arena. This mode is used for repeatable experiments.
\end{itemize}

Each simulated fighter stores position, velocity, yaw, pitch, health, attack cooldown, and control state. The local simulator is not a complete Minecraft PvP clone, but it is useful for testing planning logic quickly and consistently.

\subsection{Baseline and Standard MCTS}

The baseline non-search behavior is a distance-threshold policy: move forward when far from the opponent, sprint when very far, jump for vertical gaps, and attack when inside sword reach. This policy is useful as a rollout opponent and as a conceptual baseline because it has no explicit planning.

The standard MCTS planner clones the current duel state and searches over candidate combat actions. Each decision performs:
\begin{enumerate}
  \item \textbf{Selection}: choose a child with UCT.
  \item \textbf{Expansion}: add an untried combat action.
  \item \textbf{Simulation}: apply movement with prismarine-physics and resolve attacks when allowed.
  \item \textbf{Rollout}: continue for a short horizon using heuristic actions.
  \item \textbf{Backpropagation}: add the rollout score to all nodes on the selected path.
\end{enumerate}

The reward function values survival, opponent damage, useful spacing, cooldown timing, and attack windows. It penalizes death, large vertical mismatch, and being too far from relevant combat range.

\subsection{Creative Variations}

Four MCTS variations were implemented for comparison:
\begin{itemize}
  \item \textbf{Hybrid MCTS + heuristic evaluation}: adds a heuristic action bonus to the rollout score. The bonus rewards actions that approach at long range, attack when a valid cooldown window exists, and remain near useful melee spacing.
  \item \textbf{Risk-aware MCTS}: increases the value of safer spacing when the agent has lower health than the opponent.
  \item \textbf{High-exploration MCTS}: raises the UCT exploration constant, encouraging the search to test more candidate actions before committing.
  \item \textbf{Higher-budget MCTS}: increases iteration count and rollout depth, giving the planner more simulated futures before selecting a move.
\end{itemize}

These variations are intentionally not guaranteed to improve the result. They are designed to test different hypotheses about what matters in this combat domain: tactical heuristics, defensive survival, broader action exploration, and raw planning budget.

\section{Experimental Evaluation}

The evaluation uses physics mode because it is repeatable and fast enough for multiple runs. The report generator creates five strategy variants with eight repeated trials each. The trials use independent deterministic jitter in the starting positions and rotate through four opponent styles: balanced, defensive, aggressive, and stronger-search opponents. This avoids over-interpreting a small number of nearly identical deterministic duels. The generator also creates two stress-test scenarios: a wide-arena opening and a low-health opening.

Each logged decision records the selected action, attack flag, rollout score, visit count, distance, health difference, cooldowns, positions, and movement controls. The generated files are stored in \path{debug/report/}. The most important files are \path{aggregate-summary.csv}, \path{summary.csv}, \path{episodes.csv}, \path{action-counts.csv}, and \path{tick-series.csv}.

The measurements are winner, final health, fight length, first attack tick, number of attack decisions, average distance, and average MCTS visits. Repeated trials make it possible to report means and variances rather than relying on a single duel.

\section{Results}

Table~\ref{tab:aggregate} summarizes the repeated strategy trials. Standard MCTS and high-exploration MCTS had the best Alpha win rate at 0.625. Higher-budget MCTS won half of its runs while using the most search visits on average. The hybrid heuristic version had the lowest win rate in this setting, despite using more visits than standard MCTS. This is a useful negative result: the heuristic bonus changed behavior, but it did not automatically make the agent stronger against varied opponents.

\begin{table}[h]
\centering
\caption{Aggregate results from eight repeated trials per strategy variant.}
\label{tab:aggregate}
\small
\resizebox{\linewidth}{!}{%
\begin{tabular}{p{0.30\linewidth}rrrrrr}
\toprule
Variant & Runs & Alpha Win Rate & Mean Ticks & Var. Ticks & Mean First Attack & Mean Visits \\
\midrule
Standard MCTS & 8 & 0.625 & 56.250 & 10.188 & 13.625 & 14.304 \\
Hybrid MCTS + heuristic evaluation & 8 & 0.250 & 56.000 & 14.750 & 13.375 & 15.733 \\
Risk-aware MCTS & 8 & 0.375 & 54.625 & 21.734 & 13.375 & 14.112 \\
High-exploration MCTS & 8 & 0.625 & 55.625 & 16.734 & 13.000 & 14.024 \\
Higher-budget MCTS & 8 & 0.500 & 55.000 & 19.500 & 12.250 & 20.063 \\
\bottomrule
\end{tabular}
}
\end{table}

\begin{figure}[h]
\centering
\setlength{\fboxsep}{8pt}
\fbox{\begin{minipage}{0.86\linewidth}
\small
\textbf{Alpha win rate by strategy variant}\\[4pt]
Standard MCTS \hfill 0.625 \quad \textcolor{blue!70!black}{\rule{0.375\linewidth}{6pt}}\\[4pt]
Hybrid heuristic \hfill 0.250 \quad \textcolor{blue!70!black}{\rule{0.150\linewidth}{6pt}}\\[4pt]
Risk-aware \hfill 0.375 \quad \textcolor{blue!70!black}{\rule{0.225\linewidth}{6pt}}\\[4pt]
High exploration \hfill 0.625 \quad \textcolor{blue!70!black}{\rule{0.375\linewidth}{6pt}}\\[4pt]
Higher budget \hfill 0.500 \quad \textcolor{blue!70!black}{\rule{0.300\linewidth}{6pt}}
\end{minipage}}
\caption{Alpha win rate across repeated strategy trials. The full generated chart is also saved as \texttt{debug/report/win-rate-by-variant.svg}.}
\label{fig:winrate}
\end{figure}

\begin{figure}[h]
\centering
\setlength{\fboxsep}{8pt}
\fbox{\begin{minipage}{0.94\linewidth}
\small
\textbf{Action distribution by strategy variant}\\[6pt]
Standard MCTS\hfill
\textcolor{blue!70!black}{\rule{0.414\linewidth}{6pt}}
\textcolor{green!50!black}{\rule{0.091\linewidth}{6pt}}
\textcolor{red!70!black}{\rule{0.058\linewidth}{6pt}}
\textcolor{orange!90!black}{\rule{0.014\linewidth}{6pt}}
\textcolor{purple!80!black}{\rule{0.018\linewidth}{6pt}}
\textcolor{gray}{\rule{0.006\linewidth}{6pt}}\\[4pt]
Hybrid heuristic\hfill
\textcolor{blue!70!black}{\rule{0.417\linewidth}{6pt}}
\textcolor{green!50!black}{\rule{0.081\linewidth}{6pt}}
\textcolor{red!70!black}{\rule{0.052\linewidth}{6pt}}
\textcolor{orange!90!black}{\rule{0.028\linewidth}{6pt}}
\textcolor{purple!80!black}{\rule{0.015\linewidth}{6pt}}
\textcolor{gray}{\rule{0.005\linewidth}{6pt}}\\[4pt]
Risk-aware\hfill
\textcolor{blue!70!black}{\rule{0.396\linewidth}{6pt}}
\textcolor{green!50!black}{\rule{0.100\linewidth}{6pt}}
\textcolor{red!70!black}{\rule{0.055\linewidth}{6pt}}
\textcolor{orange!90!black}{\rule{0.015\linewidth}{6pt}}
\textcolor{purple!80!black}{\rule{0.028\linewidth}{6pt}}
\textcolor{gray}{\rule{0.006\linewidth}{6pt}}\\[4pt]
High exploration\hfill
\textcolor{blue!70!black}{\rule{0.398\linewidth}{6pt}}
\textcolor{green!50!black}{\rule{0.107\linewidth}{6pt}}
\textcolor{red!70!black}{\rule{0.049\linewidth}{6pt}}
\textcolor{orange!90!black}{\rule{0.018\linewidth}{6pt}}
\textcolor{purple!80!black}{\rule{0.018\linewidth}{6pt}}
\textcolor{gray}{\rule{0.009\linewidth}{6pt}}\\[4pt]
Higher budget\hfill
\textcolor{blue!70!black}{\rule{0.405\linewidth}{6pt}}
\textcolor{green!50!black}{\rule{0.095\linewidth}{6pt}}
\textcolor{red!70!black}{\rule{0.057\linewidth}{6pt}}
\textcolor{orange!90!black}{\rule{0.017\linewidth}{6pt}}
\textcolor{purple!80!black}{\rule{0.025\linewidth}{6pt}}
\textcolor{gray}{\rule{0.001\linewidth}{6pt}}\\[8pt]
\textcolor{blue!70!black}{\rule{0.04\linewidth}{6pt}} sprint-advance \quad
\textcolor{green!50!black}{\rule{0.04\linewidth}{6pt}} advance \quad
\textcolor{red!70!black}{\rule{0.04\linewidth}{6pt}} close-attack \quad
\textcolor{orange!90!black}{\rule{0.04\linewidth}{6pt}} jump-sprint \quad
\textcolor{purple!80!black}{\rule{0.04\linewidth}{6pt}} jump-advance \quad
\textcolor{gray}{\rule{0.04\linewidth}{6pt}} hold/other
\end{minipage}}
\caption{Stacked action distribution by strategy variant. Bar segments are proportional to the percentage of logged decisions for that variant. The full generated chart is saved as \texttt{debug/report/action-distribution-by-variant.svg}.}
\label{fig:actiondistribution}
\end{figure}

Table~\ref{tab:examples} shows individual stress-test runs. The wide-arena example delays the first attack until tick 20 because the fighters start farther apart. The low-health example is won by Bravo, which shows that the simulator is not hard-coded to make Alpha win.

\begin{table}[h]
\centering
\caption{Stress-test examples from the generated summary.}
\label{tab:examples}
\small
\begin{tabular}{p{0.28\linewidth}lrrrrr}
\toprule
Scenario & Winner & Episodes & Ticks & Attacks & Alpha HP & First Attack \\
\midrule
Wide arena example & Alpha & 107 & 54 & 8 & 8 & 20 \\
Low-health example & Bravo & 70 & 35 & 6 & 0 & 13 \\
\bottomrule
\end{tabular}
\end{table}

Across all runs, movement actions dominate early decisions. Figure~\ref{fig:actiondistribution} shows that every strategy uses \texttt{sprint-advance} for roughly two thirds of logged decisions, which is expected because most duels begin outside melee range. The variants still differ in smaller ways: high-exploration MCTS uses the most \texttt{advance}, risk-aware MCTS uses more \texttt{jump-advance}, and hybrid MCTS uses more \texttt{jump-sprint}. Attack actions appear only after the fighters have moved into sword reach and cooldown conditions allow damage. This supports the claim that the planner is coordinating approach and attack timing rather than simply attacking every tick.

\section{Discussion}

The results suggest that exploration and opponent diversity matter. Standard MCTS and high-exploration MCTS achieved the best Alpha win rate, while higher-budget MCTS produced the largest average visit count and the earliest mean first attack tick. More search effort therefore changed behavior, but it did not dominate every opponent style.

The hybrid heuristic variation is more subtle. It increased average visits relative to standard MCTS but achieved a lower win rate. A likely explanation is that the heuristic bonus made the agent overvalue immediate attack windows or close spacing, while defensive and stronger-search opponents punished that positioning. This is a good example of reward shaping risk: a plausible hand-written heuristic can produce worse strategic behavior if it is not aligned with long-term survival.

The risk-aware variation won three of eight repeated trials. In the low-health stress test, however, Alpha still lost. This indicates that defensive spacing alone is not enough when the health and damage disadvantage is large. More realistic defensive behavior would need knockback, shield timing, terrain use, or a better opponent model.

The main limitation is simulator fidelity. The local prismarine-physics duel approximates movement and cooldowns, but real Minecraft combat includes latency, knockback, shields, item effects, terrain, server version differences, and human unpredictability. Another limitation is sample size: eight repeated trials per variant is enough to study mean and variance at a basic level, but not enough for strong statistical claims. A larger study should run dozens or hundreds of randomized openings per variant.

\section{Conclusion}

This project implements a Minecraft sword PvP agent based on Monte Carlo Tree Search and evaluates several creative variations of the algorithm. The final system supports Mineflayer server play and local prismarine-physics experiments. The report data includes 42 duel runs and 4,598 logged decisions.

The experiments show that MCTS can produce plausible combat behavior in this environment: the agent closes distance, waits until melee range, and attacks in cooldown-limited bursts. The variations also show that changing the search procedure changes outcomes. In this experiment, standard MCTS and high-exploration MCTS achieved the best win rate, while higher-budget MCTS attacked earliest and used the most search effort. The hybrid heuristic bonus changed behavior without improving win rate. Future work should use larger randomized batches, stronger opponent models, and a more faithful PvP simulator with knockback, shields, terrain, and item differences.

\section*{Acknowledgements}

No external human collaboration was used for the implementation unless otherwise stated by the author. AI assistance was used to help inspect the brief, improve code structure, generate experiment material, and improve the readability of the report. The project uses open-source Minecraft bot libraries including Mineflayer, prismarine-physics, minecraft-data, prismarine-block, and vec3.

\section*{AI and Implementation Statement Summary}

A separate \texttt{statement.pdf} is included with the submission. It explains which parts of the work were based on existing libraries, which parts were created or edited with AI assistance, and which parts represent the student's own implementation and design work.

\begin{thebibliography}{9}

\bibitem{coulom2006}
R. Coulom.
\newblock Efficient selectivity and backup operators in Monte-Carlo tree search.
\newblock In \emph{Computers and Games}, Lecture Notes in Computer Science, vol. 4630, pp. 72--83, 2007.
\newblock \url{https://doi.org/10.1007/978-3-540-75538-8_7}.

\bibitem{kocsis2006}
L. Kocsis and C. Szepesvari.
\newblock Bandit based Monte-Carlo planning.
\newblock In \emph{Machine Learning: ECML 2006}, Lecture Notes in Computer Science, vol. 4212, pp. 282--293, 2006.
\newblock \url{https://doi.org/10.1007/11871842_29}.

\bibitem{browne2012}
C. B. Browne, E. Powley, D. Whitehouse, S. M. Lucas, P. I. Cowling, P. Rohlfshagen, S. Tavener, D. Perez, S. Samothrakis, and S. Colton.
\newblock A survey of Monte Carlo tree search methods.
\newblock \emph{IEEE Transactions on Computational Intelligence and AI in Games}, 4(1):1--43, 2012.
\newblock \url{https://doi.org/10.1109/TCIAIG.2012.2186810}.

\bibitem{silver2016}
D. Silver, A. Huang, C. J. Maddison, A. Guez, L. Sifre, G. van den Driessche, J. Schrittwieser, I. Antonoglou, V. Panneershelvam, M. Lanctot, and others.
\newblock Mastering the game of Go with deep neural networks and tree search.
\newblock \emph{Nature}, 529:484--489, 2016.
\newblock \url{https://doi.org/10.1038/nature16961}.

\bibitem{mineflayer}
Mineflayer contributors.
\newblock Mineflayer: Create Minecraft bots with a powerful, stable, and high level JavaScript API.
\newblock GitHub repository, PrismarineJS.
\newblock \url{https://github.com/PrismarineJS/mineflayer}.

\bibitem{prismarinephysics}
PrismarineJS contributors.
\newblock prismarine-physics.
\newblock GitHub repository, PrismarineJS.
\newblock \url{https://github.com/PrismarineJS/prismarine-physics}.

\end{thebibliography}

\end{document}
